% Section 6 -- Experimental Setup
% Needs: \usepackage{amsmath, booktabs}
% Hardware and job details come from the Slurm job logs. Run counts come from results/.

\section{Experimental Setup}
\label{sec:setup}

All GPU runs and CPU timing runs ran as Slurm batch jobs on a shared university cluster.
Each node has two NVIDIA A40 GPUs with 48\,GB each, two 16-core CPUs and 192\,GB of main memory.
We used Python 3.11, CUDA-Q 0.16 and SciPy 1.17.
Each GPU job used one GPU, and every job had a 24-hour limit.
Table~\ref{tab:experiments} lists the settings of each experiment, all with $\lambda = 0.5$.
Above $n = 24$ we know only one optimal schedule and its mirror image, so $P(\text{opt})$ at those sizes may be slightly too low.

\begin{table}[h]
  \centering
  \caption{Settings of each experiment. Instances are counted per size. Sizes go up in steps of 2 except in E3, which uses every size. In E2, INTERP also runs depths 1, 3 and 5.}
  \label{tab:experiments}
  \begin{tabular}{@{}llllll@{}}
    \toprule
    Exp. & Purpose & Simulator & Sizes $n$ & Instances & Depth $p$ \\
    \midrule
    E1 & Depth sweep & CPU fp64 & 8--12 & 10 & 1--6 \\
       &             & GPU fp32 & 14--20 & 10 & 1--6 \\
    E2 & Starting angles & GPU fp32 & 10--16 & 10 & 2, 4, 6 \\
    E3 & Classical methods & none & 8--20 & 10 & none \\
       &                   & none & 21--28 & 1 & none \\
    E4 & Timing & CPU fp64 & 10--20 & 3 & 3 \\
       &        & GPU fp32 and fp64 & 10--30 & 3 & 3 \\
    E5 & Reused angles at large $n$ & GPU fp32 & 20--28 & 10 & 1--6 \\
    E6 & Gate noise & CPU density matrix & 6, 8 & 10 & 1--4 \\
    \bottomrule
  \end{tabular}
\end{table}

E4 changes only the target between runs.
Runs stop after different numbers of evaluations, so we compare the time per evaluation.
The CPU runs used 8 threads rather than the whole node.
Since qpp-cpu supports only fp64, the fair comparison is CPU fp64 against GPU fp64.
The nodes were shared, so timings reflect normal working conditions.
Two settings hit the 24-hour limit, which leaves one of three instances for GPU fp64 at $n = 30$ and two of three for the CPU at $n = 20$.

Because E1 mixes CPU fp64 and GPU fp32 results, we also ran $n = 12$ on the GPU.
The approximation ratio changed by less than 0.001 in 58 of 60 runs and the best sampled cost never changed, so fp32 does not affect our conclusions.

Precision also limits the largest problem a GPU can hold.
The simulator stores all $2^n$ amplitudes of the state, at 8 bytes each in fp32 and 16 bytes in fp64.
Each extra qubit doubles the memory, so a 48\,GB A40 holds at most 32 qubits in fp32 and 31 in fp64.
